\section{Prototypische Umsetzung}

Der Prototyp setzt das beschriebene Verfahren als Python-Konsolenanwendung um.
Das Projekt trennt Datenbeschaffung, Optimierung und gemeinsame Hilfsfunktionen
in die Verzeichnisse \texttt{api}, \texttt{optimizer} und \texttt{helpers} unter
\texttt{src}. \texttt{main.py} verbindet diese Teile und gibt die empfohlenen
Reisepläne aus; \texttt{benchmark.py} ermöglicht Vergleichsläufe auf denselben
gespeicherten Angeboten. Während der Optimierung erfolgen keine Anbieterabfragen.

Die Problemvorgaben werden zentral in \texttt{config.py} festgelegt: Startort
und mögliche Abflughäfen, Reisezeitfenster und Reisedauer, Kandidaten- und
Pflichtziele, Anzahl der besuchten Städte sowie Personenzahl. Hinzu kommen
Mindestaufenthalte, nutzbare Tagesstunden und Obergrenzen für Budget, Emissionen
und die Dauer einzelner Verbindungen. Die Unterkunftsplanung lässt sich über
\texttt{OPTIMIZE\_HOTELS} ein- und ausschalten; bei deaktivierter Planung
entfallen Hotelauswahl und sämtliche Unterkunftsbeiträge zur Bewertung,
nicht jedoch die Aufenthaltsbedingungen. Bei aktivierter Planung begrenzen
Sterneklasse, Gästebewertung, Bewertungsanzahl und Nachtpreis die Hotelauswahl.
Suchparameter und Bewertungsgewichte stehen getrennt in \texttt{settings.toml},
Zeitreserven und Ortsinformationen in \texttt{assumptions.toml}.

Die abstrakte Repräsentation wird als Python-Datenklasse mit Listen für
Route, Nächte und Auswahlgene umgesetzt. Verkehrsmittel und Abfahrtszeitfenster
werden getrennt gespeichert; das Fenster wird aus einer Stundenangabe
abgeleitet. Zusätzlich wird der Abreisetag innerhalb des vorgegebenen
Reisefensters mitoptimiert und als Tagesoffset gespeichert. Anders als beim
vereinfachten Nachtversatz des Modells ergeben sich die Ankunftstage aus
konkreten, zeitzonenbezogenen Fahrplanzeiten einschließlich Zu- und
Abgangsreserven. Auch der Diskomfort wird verfeinert: Modus, Reisezeit,
Schlafunterbrechungen, lange Fahrten und Wartezeiten am Hotel gehen ein.

Der Decoder folgt den Transportpräferenzen ohne Rückwärtssuche; fehlt eine
Fortsetzung, bleibt der Kandidat unzulässig. Bei nicht erfüllbaren Präferenzen
hat zunächst das Verkehrsmittel, danach das Zeitfenster Vorrang.
Abweichend von einer pauschalen Dominanzreduktion innerhalb einer Klasse
werden Transportangebote nur bei übereinstimmenden Zeit- und
Verbindungseigenschaften verglichen, da andere Ankunftszeiten spätere
Aufenthalte verändern können. Bei Hotels werden zusätzlich Check-in- und
Check-out-Zeiten berücksichtigt. Zwischenspeicher ersetzen eine allgemeine
inkrementelle Neubewertung einzelner Etappen. Die im Modell vorgesehene
Flugbeschränkung nach Entfernung oder Zeitvorteil ist noch nicht umgesetzt;
die relative Budget- und Emissionsstrafe verwendet jeweils Gewicht eins,
ergänzt um die Anzahl struktureller Verstöße.

Flug- und Hotelangebote stammen aus Schnittstellen zu Google Flights und
Google Hotels, Bodenverbindungen aus Transitous sowie FlixBus und FlixTrain.
Der Bahn- und Buszugriff ist im Prototyp auf unterstützte europäische Städte
und Strecken bis 1.350\,km Luftlinie begrenzt.
Wo verbindliche Fahrpreise fehlen, werden sie ersatzweise geschätzt:
Der Transitous-Bahnprovider verwendet einen entfernungsabhängigen Näherungstarif,
während Flix-Angebotspreise übernommen werden.

Mangels verbindungsspezifischer Emissionsangaben werden die Bodenemissionen
aus Entfernung und einem Faktor je Personenkilometer angenähert, entsprechend
dem gebräuchlichen Bezug von Verkehrsemissionen auf die Verkehrsleistung
\parencite[Vgl.][]{ubaEmissionsdaten}:
\begin{equation}
  \widehat e(v)
  = N_{\mathrm{Reisende}}\,
    \underbrace{1{,}25\,s_{\mathrm{Luft}}(v)}_{\text{geschätzte Strecke in km}}\,
    \frac{r(v)}{1000},
  \qquad
  r(v)=
  \begin{cases}
    14 & \text{Bahn},\\
    27 & \text{Bus oder Bahn+Bus}.
  \end{cases}
\end{equation}
Dabei ist $r$ in Gramm je Personenkilometer angegeben, das Ergebnis in
Kilogramm. Der Umwegfaktor und die beiden Zahlenwerte sind vereinfachende
Prototypannahmen, keine aus der zitierten Quelle übernommenen aktuellen
Emissionsfaktoren. Für Hotels werden pauschale Nachtwerte nach Unterkunftsart
und Sterneklasse mit der Übernachtungszahl multipliziert.
Flugemissionen werden dagegen aus Google Flights übernommen; Angebote ohne
Emissionswert werden ausgeschlossen, nicht nachgeschätzt.

Die Ergebnisse sind Planungsvorschläge, keine buchbaren Gesamtangebote.
Gespeicherte Preise können veralten; die Summe einzelner Hotelnachtangebote
garantiert keine durchgehende Verfügbarkeit zum ausgewiesenen Gesamtpreis.
Flughafenzubringer und sonstige lokale Wege werden nur über Zeitreserven
berücksichtigt, nicht mit eigenen Kosten oder Emissionen; Aktivitäten und
optionale Zusatzleistungen bleiben ebenfalls außerhalb der Bewertung.
Die Schätzwerte dienen damit dem Vergleich von Alternativen, nicht einer
vollständigen Kosten- oder Klimabilanz der tatsächlichen Reise.
